A GLev ciphertext is a list of GLWE ciphertexts that encrypt the list of plaintexts $\dfrac{q}{\beta^1}M, \dfrac{q}{\beta^2}M, \ldots, \dfrac{q}{\beta^l}M$, where $M$ is a plaintext encoded in a polynomial. Note that each $i$-th GLWE ciphertext of a GLev ciphertext uses a different plaintext scaling factor, which is: $\Delta_i = \dfrac{q}{\beta^i}$. The structure of GLev ciphertext is visually depicted in \autoref{fig:glev}.


For these integer-ring weights require an integer $\beta\ge2$ and $\beta^l\mid q$. The inequality $t\le\beta$ is needed only if the first row is intended to recover coefficients represented in $\{0,\ldots,t-1\}$. GLev is also used as a gadget ciphertext for a general ring message; in that use, each row need not separately decode the original message. Exact reconstruction with the descending gadget for every residue additionally requires $q=\beta^l$, as explained in \autoref{sec:decomp}.


\subsection{Encryption}
\label{subsec:glev-enc}

\begin{tcolorbox}[title={\textbf{\tboxlabel{\ref*{subsec:glev-enc}} GLev Encryption}}]


$\textsf{GLev}_{S, \sigma}^{\beta, l}(M) = \Bigl \{\textsf{GLWE}_{S, \sigma}\left(\dfrac{q}{\beta^i} M + E_i\right)  \Bigr \}_{i=1}^{l} \in \mathcal{R}_{\langle n, q \rangle }^{(k+1) \cdot l}$
\end{tcolorbox}

\begin{figure}[h!]
    \centering
  \includegraphics[width=1.0\linewidth]{figures/TFHE-fig2.pdf}
  \caption{An illustration of a GLev ciphertext }
  \label{fig:glev}
\end{figure}



\subsection{Decryption}

The first row has phase $(q/\beta)M+E_1\pmod q$. If $\|E_1\|_\infty<q/(2\beta)$, nearest-integer decoding recovers $M\bmod\beta$ coefficientwise. It recovers a chosen original representative only when that representative is uniquely determined by its class modulo $\beta$, for example coefficients in $\{0,\ldots,t-1\}$ with $t\le\beta$. Row $i$ similarly decodes modulo $\beta^i$ if $\|E_i\|_\infty<q/(2\beta^i)$. Later rows have a smaller error tolerance; they are not necessarily undecodable, nor is separate decoding required for gadget operations.

\subsection{Lev and RLev}

Lev is GLev with $n=1$. RLev is GLev with $k=1$.

